\section{Appendix: Seasonal Votive Texts}
These alternatives are not used on 25 October 2026. After Septuagesima the Tract replaces the Alleluia while the Gradual remains. In Eastertide the two-verse Alleluia replaces the ordinary Gradual and Alleluia. The printed feast rubric governs these alternatives; they do not supply additional readings for the dated feast.
\subsubsection{Tract: inactive seasonal alternative}
\properlocus{1962 Missal, p. 698, no. 3943; Lasance, p. 1259.}
\begin{latinproper}
Ipse invocábit me: Pater meus es tu: Deus meus, et suscéptor salútis meæ. \textnormal{V.} Et ego primogénitum ponam illum, excélsum præ régibus terræ. \textnormal{V.} Et ponam in sǽculum sǽculi semen eius: et thronum eius sicut dies cæli.
\end{latinproper}
\begin{englishwitness}{Lasance--Walsh}
He shall cry out to Me: Thou art My Father: My God, and the support of My salvation. \textnormal{V.} And I will make Him My firstborn, high above the kings of the earth. \textnormal{V.} And I will make His seed to endure for evermore: and His throne as the days of heaven.
\end{englishwitness}
\subsubsection{Paschal Alleluia: inactive seasonal alternative}
\properlocus{1962 Missal, p. 698, no. 3944; Lasance, p. 1259.}
\begin{latinproper}
Allelúia, allelúia. \textnormal{V.} \textnormal{Dan. 7, 14} Potéstas eius, potéstas ætérna, quæ non auferétur: et regnum eius, quod non corrumpétur. Allelúia. \textnormal{V.} \textnormal{Apoc. 19, 16} Habet in vestiménto et in fémore suo scriptum: Rex regum et Dóminus dominántium. Allelúia.
\end{latinproper}
\begin{englishwitness}{Lasance--Walsh}
Alleluia, alleluia. \textnormal{V.} \textnormal{Dan. 7, 14.} His power is an everlasting power that shall not be taken away: and His kingdom that shall not be destroyed, Alleluia. \textnormal{V.} \textnormal{Apoc. 19, 16.} He hath on His garment, and on His thigh written: King of kings, and Lord of lords, Alleluia.
\end{englishwitness}
